\documentclass[11pt]{amsart}
\usepackage[margin=1.1in]{geometry}
\usepackage{amsmath,amssymb,amsthm}
\usepackage{enumitem}
\usepackage[hidelinks]{hyperref}

\newtheorem{theorem}{Theorem}
\newtheorem{lemma}[theorem]{Lemma}
\newtheorem{corollary}[theorem]{Corollary}
\theoremstyle{remark}
\newtheorem{remark}[theorem]{Remark}

\newcommand{\F}{\mathbb{F}_2}
\newcommand{\W}{W}
\newcommand{\nl}{\mathrm{nl}}
\newcommand{\sg}{\sigma}
\newcommand{\wt}{\mathrm{wt}}
\newcommand{\Rad}{\mathrm{Rad}}

\title[Balanced functions in even dimension]{A sharper upper bound on the nonlinearity\\ of balanced Boolean functions in even dimension}
\author{Daniel Ayomide Shittu}
\address{Rowan University, Glassboro, NJ, USA}
\date{November 2026}

\begin{document}
\begin{abstract}
Let $f$ be a balanced Boolean function in $n=2m$ variables. The classical bound
is $\nl(f)\le 2^{n-1}-2^{m-1}-2$. We prove that for every $m\ge 4$,
\[
  \nl(f)\;\le\;2^{\,n-1}-2^{\,m-1}-4 .
\]
This answers Open Problem~4 of Lou, Xu and Wang in the affirmative. The bound is
attained for $n=8$ and $n=10$, so it recovers $\nl_{\mathrm{bal}}(8)=116$ and
$\nl_{\mathrm{bal}}(10)=492$ with a single, computer-free argument. For $n=12,14,16$ it
gives $2012$, $8124$ and $32636$. The proof uses only the Walsh transform modulo~16,
Parseval's identity and the spectrum of quadratic forms. Every step has been checked
in the Lean~4 proof assistant, in two self-contained files without external
libraries.
\end{abstract}

\maketitle

\section{Introduction}

For $f:\F^n\to\F$ write $\sg(b)=(-1)^b$ and
\[
  \W_f(a)=\sum_{x\in\F^n}\sg\bigl(f(x)\oplus a\cdot x\bigr),\qquad
  \nl(f)=2^{n-1}-\tfrac12\max_a|\W_f(a)| .
\]
The function $f$ is \emph{balanced} if it takes the value $1$ exactly $2^{n-1}$ times,
that is, $\W_f(0)=0$. Let $\nl_{\mathrm{bal}}(n)$ be the largest nonlinearity of a balanced
function in $n$ variables. For even $n=2m$, bent functions reach $2^{n-1}-2^{m-1}$ but are
never balanced, and the nonlinearity of a balanced function is even. This gives the
classical bound $\nl_{\mathrm{bal}}(2m)\le 2^{2m-1}-2^{m-1}-2$ \cite{SZZ93,Carlet21}.
Dobbertin's construction \cite{Dob94} gives
$\nl_{\mathrm{bal}}(2m)\ge 2^{2m-1}-2^m+\nl_{\mathrm{bal}}(m)$, and he conjectured that this is
an equality.

In 2026 the cases $n=8$ and $n=10$ were settled: $\nl_{\mathrm{bal}}(8)=116$ by Yuan, Hu, Zhang, Zhang and Wu \cite{YHZZW26}
(with computer assistance) and by Lou, Xu and Wang \cite{LXW26} (by hand), and
$\nl_{\mathrm{bal}}(10)=492$ by Xu and Wang \cite{XW26}. Lou, Xu and Wang asked
\cite[Open Problem~4]{LXW26}: \emph{if $f$ is balanced on $2n$ variables, is
$\nl(f)\le 2^{2n-1}-2^{n-1}-4$ for every $n\ge 4$?} We answer this question.

\begin{theorem}\label{thm:main}
Let $m\ge 4$ and let $f$ be a balanced Boolean function in $n=2m$ variables. Then
\[
  \nl(f)\le 2^{\,2m-1}-2^{\,m-1}-4 .
\]
\end{theorem}

\begin{corollary}\label{cor}
$\nl_{\mathrm{bal}}(8)=116$, $\nl_{\mathrm{bal}}(10)=492$, $\nl_{\mathrm{bal}}(12)\le 2012$,
$\nl_{\mathrm{bal}}(14)\le 8124$ and $\nl_{\mathrm{bal}}(16)\le 32636$. The same bounds hold for
every component of a balanced $(2m,r)$ S-box, and in particular for permutations of
$\F^{2m}$, $m\ge4$.
\end{corollary}

The values at $n=8,10$ are attained by Dobbertin's construction. At $n=12,14,16$ that
construction gives $2010$, $8120$ and $32628$, so the remaining gaps are $2$, $4$ and $8$.
The hypothesis $m\ge 4$ is needed: $\nl_{\mathrm{bal}}(6)=26$ exceeds the value $24$ given by the formula
(the proof uses $16\mid 2^m$, see Step~3).

\medskip\noindent\textbf{Relation to earlier work.}
The new content is the general statement for all $m\ge 4$. For $n=8$ and $n=10$ the
results are due to \cite{YHZZW26,LXW26,XW26}; we give independent proofs of them, inside
the same argument. Our starting point, a congruence for $\W_f$ modulo~16 that produces a
quadratic form, is also the starting point of \cite{LXW26}. As far as we can tell from the
abstracts of \cite{XW26,YHZZW26} and the text of \cite{LXW26}, none of them states a bound
for general even $n$ for single Boolean functions.

\medskip\noindent\textbf{Formal proof.}
Theorem~\ref{thm:main} is proved in Lean~4 (Section~\ref{sec:lean}). The files use no
library beyond Lean's core, contain no \texttt{sorry} and no added axioms, and depend
only on Lean's three standard axioms.

\medskip\noindent\textbf{Plan.} Section~\ref{sec:prelim} collects three standard facts.
Section~\ref{sec:main} proves the theorem for $n\ge 12$ and reduces $n=8,10$ to a short list
of parameters. Sections~\ref{sec:ten} and~\ref{sec:eight} finish $n=10$ and $n=8$.

\section{Preliminaries}\label{sec:prelim}

For $g:\F^N\to\mathbb{Z}$ put $\W_g(a)=\sum_x g(x)\sg(a\cdot x)$, so
$\W_f=\W_{\sg\circ f}$. We use Parseval, $\sum_a \W_g(a)^2=2^N\sum_x g(x)^2$, the
adjoint identity $\sum_a \W_g(a)h(a)=\sum_x g(x)\W_h(x)$, and inversion
$\W_{\W_g}=2^N g$.

\begin{lemma}[Walsh values modulo 16]\label{lem:mod16}
Let $f$ be balanced on $n\ge 4$ variables, $v=\bigoplus_{f(x)=1}x$, and
$N(a)=\#\{x: f(x)=1,\ a\cdot x=1\}$. Then $\W_f(a)=4N(a)-2^n$ for $a\neq 0$, and
\[
  \W_f(a)\equiv 4\,[a\cdot v]+8\,Q(a)\pmod{16},\qquad Q(a)=\tbinom{N(a)}{2}\bmod 2 .
\]
The function $Q$ is quadratic and $Q(0)=0$. In particular $4\mid \W_f$ everywhere.
\end{lemma}

\begin{proof}
Counting the four classes of $x$ according to $f(x)$ and $a\cdot x$ gives
$\W_f(a)=4N(a)-2^n$. Modulo $16$ this is $4(N\bmod 4)$, and $N\bmod 2=a\cdot v$ while
$\binom N2\bmod 2=1$ exactly when $N\equiv 2,3\pmod 4$. Finally
$\binom{N(a)}2=\sum_{\{x,y\}}(a\cdot x)(a\cdot y)$, summed over pairs in the support of
$f$, is a polynomial of degree at most $2$ in $a$.
\end{proof}

For a quadratic $q:\F^N\to\F$, its radical is
$\Rad(q)=\{u: q(a\oplus u)\oplus q(a)\text{ is constant in }a\}$.

\begin{lemma}[quadratic spectra]\label{lem:quad}
Let $q$ be quadratic on $N=2k+1$ variables. There is $r\le k$ with $|\Rad(q)|=2^{N-2r}$,
$\W_q(x)\in\{0,\pm 2^{N-r}\}$ for all $x$, and $\sum_x|\W_q(x)|=2^{N+r}$.
\end{lemma}

This is the standard description of quadratic forms \cite[Ch.~5]{Carlet21}.

\section{The general argument}\label{sec:main}

Fix $k\ge 3$ and put $n=2k+2$, $B=2^k$. So $2^n=4B^2$, $2^{m}=2B$ and $16\mid 2B$.
By Lemma~\ref{lem:mod16}, $\nl(f)\ge 2^{n-1}-2^{m-1}-3$ is equivalent, for balanced $f$, to
$|\W_f|\le 2B+4$ everywhere. Theorem~\ref{thm:main} therefore follows from:

\begin{theorem}\label{thm:walsh}
For $k\ge 3$ there is no balanced $f$ on $2k+2$ variables with $|\W_f(a)|\le 2B+4$ for all $a$.
\end{theorem}

Assume such an $f$ exists. We derive a contradiction in nine steps. Throughout, $v$ and $Q$
are as in Lemma~\ref{lem:mod16}.

\smallskip\noindent\textbf{Step 1 ($v\neq0$).}
If $v=0$, then $8\mid \W_f$, so $|\W_f|\le 2B$ (because $8\mid 2B$). Since $\W_f(0)=0$,
Parseval gives $2^{2n}=\sum\W_f^2\le(2^n-1)4B^2<2^{2n}$, a contradiction. After a linear change of
variables (which preserves balance and permutes the spectrum) we may assume $a\cdot v=a_0$.
Write $a=(a_0,a')$ with $a'\in\F^{2k+1}$, $H_0=\{a_0=0\}$, $H_1=\{a_0=1\}$, and
$Q_c(a')=Q(c,a')$.

\smallskip\noindent\textbf{Step 2 (the two halves).}
Define $S,D:\F^{2k+1}\to\{-1,0,1\}$ by $S(x')=\sg(f(0,x'))$ if $f(0,x')=f(1,x')$ and $0$
otherwise, and $D(x')=\sg(f(0,x'))$ if $f(0,x')\ne f(1,x')$ and $0$ otherwise. Then
\[
  \W_f(0,a')=2\W_S(a'),\qquad \W_f(1,a')=2\W_D(a').
\]
The supports of $S$ and $D$ partition $\F^{2k+1}$. Write $|\mathrm{supp}\,D|=B^2+2t$; here $t$ is
an integer, since $\W_D(0)=\tfrac12\W_f(1,0)\equiv 2\pmod 4$ makes $|\mathrm{supp}\,D|$ even.
Parseval on each half gives
\[
  \sum_{a'}\W_f(0,a')^2=8B^4-16B^2t,\qquad \sum_{a'}\W_f(1,a')^2=8B^4+16B^2t .
\]

\smallskip\noindent\textbf{Step 3 (the $H_0$ bound).}
On $H_0$ we have $\W_f\equiv 8Q_0\pmod{16}$. Since $16\mid 2B$, this gives $|\W_f|\le 2B$ where
$Q_0=0$ and $|\W_f|\le 2B-8$ where $Q_0=1$; also $\W_f(0)=0$. Comparing with Step~2,
\begin{equation}\label{eq:H0}
  8(B-2)\,\wt(Q_0)\;\le\;B^2(4t-1).
\end{equation}
In particular $t\ge1$.

\smallskip\noindent\textbf{Step 4 (the $H_1$ bound).}
On $H_1$ put $s(a')=\W_f(1,a')\,\sg(Q_1(a'))$. By Lemma~\ref{lem:mod16}, $s\equiv 4\pmod{16}$,
and $|s|\le 2B+4$, so $-(2B-4)\le s\le 2B+4$. Hence $(2B+4-s)(2B-4+s)\ge0$, that is,
$s^2\le 4B^2-16+8s$. Summing over $a'$ and using Step~2,
\begin{equation}\label{eq:L}
  L:=\sum_{a'}s(a')\;\ge\;2B^2(t+2),
\end{equation}
with equality only if $s(a')\in\{2B+4,\,-(2B-4)\}$ for every $a'$. By the adjoint identity,
$L=2\sum_{x'}D(x')\,\W_{Q_1}(x')$.

\smallskip\noindent\textbf{Step 5 (the rank is maximal).}
Since $Q$ is quadratic, $Q_0\oplus Q_1$ is affine; so $Q_0,Q_1$ are quadratic with the same
radical. Apply Lemma~\ref{lem:quad} to $Q_1$ ($N=2k+1$) with rank $r$. Since $|D|\le1$,
$L\le 2\sum|\W_{Q_1}|=4B^2\cdot 2^r$, and \eqref{eq:L} gives $2^{r+1}\ge t+2$. If $r=0$
this contradicts $t\ge1$. If $r\ge1$, then $|\W_{Q_0}(0)|\le 2^{2k+1-r}\le B^2$, so
$\wt(Q_0)\ge B^2/2$; then \eqref{eq:H0} gives $t\ge B-1$ and so $2^{r+1}\ge B+1$, i.e.\ $r=k$.
Thus $|\Rad(Q_1)|=|\Rad(Q_0)|=2$. A linear change of the coordinates $a'$ (which keeps
$a\cdot v=a_0$) makes $\Rad=\{0,e\}$ with $e$ the first unit vector of $\F^{2k+1}$. We write
$a'=(a_1,a'')$ with $a''\in\F^{2k}$.

\smallskip\noindent\textbf{Step 6 (a lower bound on $t$).}
By Lemma~\ref{lem:quad} with $r=k$, $\W_{Q_0}(0)\in\{0,\pm2B\}$, so $\wt(Q_0)\in\{B^2,B^2\pm B\}$.
Putting $\wt(Q_0)\ge B^2-B$ into \eqref{eq:H0} gives
\begin{equation}\label{eq:tlow}
  t\ge 2B-5 .
\end{equation}

\smallskip\noindent\textbf{Step 7 (a bent function inside $Q_1$).}
Since $e\in\Rad(Q_1)$, $Q_1(a\oplus e)=Q_1(a)\oplus\varepsilon$ for a constant $\varepsilon$.
Pairing $a$ with $a\oplus e$ shows that $\W_{Q_1}(x_1,x'')=0$ unless $x_1=\varepsilon$, and
Lemma~\ref{lem:quad} with Parseval then forces $|\W_{Q_1}(\varepsilon,x'')|=2B$ for all
$x''$. So $\W_{Q_1}(\varepsilon,x'')=2B\,\sg(b(x''))$ for a function $b$ on $2k$ variables;
$q(a'')=Q_1(0,a'')$ is bent with dual $b$. Put $\tau(a'')=\sg(q(a''))$; then
$\W_{\sg\circ b}=B\tau$ and $\W_\tau=B\,\sg\circ b$.

\smallskip\noindent\textbf{Step 8 (counting the defects).}
Let $D^*(x'')=D(\varepsilon,x'')$ and $\bar D(x'')=D(1\oplus\varepsilon,x'')$, and write
$D^*=\sg\circ b+E^*$. Let
\[
  m_1=\#\{D^*=0\},\qquad m_2=\#\{D^*=-\sg\circ b\},\qquad \delta=m_1+2m_2 .
\]
Then $\sum D^*\,\sg(b)=B^2-\delta$, so by Step~7, $L=4B(B^2-\delta)$, and \eqref{eq:L} becomes
\begin{equation}\label{eq:delta}
  2\delta\le 2B^2-2B-Bt .
\end{equation}

\smallskip\noindent\textbf{Step 9 (a pointwise inequality and the main estimate).}
Put $u=\tau\,\W_{E^*}$ and $z=\tau\,\W_{\bar D}$ (functions of $a''$). Splitting $\W_D$ by
the first coordinate and using $\W_{D^*}=B\tau+\W_{E^*}$,
\[
  \tau\,\W_f(1,0,a'')=2(B+u+z),\qquad \sg(\varepsilon)\,\tau\,\W_f(1,1,a'')=2(B+u-z).
\]
Since $|\W_f|\le 2B+4$, $|B+u\pm z|\le B+2$; hence $|z|\le 2-u$ and $z^2\le(2-u)^2$. Now sum
over $a''$. Parseval and Step~7 give
\[
  \textstyle\sum z^2=B^2(2t+m_1),\quad \sum u=\sum_{x''}E^*\,\W_\tau=-B\delta,\quad
  \sum u^2=B^2(m_1+4m_2),
\]
using $|\mathrm{supp}\,\bar D|=(B^2+2t)-(B^2-m_1)$ and $E^*\sg(b)\in\{0,-1,-2\}$. Hence
\begin{equation}\label{eq:star}
  2Bt\le 4B+4\delta+4Bm_2 .\tag{$\star$}
\end{equation}
With $2m_2\le\delta$ and \eqref{eq:delta}, this gives $(B+4)\,t\le 2B^2+2B$. Together with
\eqref{eq:tlow}, $2B^2+3B-20\le 2B^2+2B$, so $B\le 20$.

\smallskip
This proves Theorem~\ref{thm:walsh} for $k\ge5$, i.e.\ $n\ge 12$. For $B=16$ we get
$t=27$. For $B=8$ we get $t\in\{11,12\}$. These two cases are finished below.

\section{The case $n=10$}\label{sec:ten}

Here $B=16$, $t=27$, and \eqref{eq:delta} reads $\delta\le 24$.

\emph{(a) $Q_0$ is invariant under $e$.} By \eqref{eq:H0}, $112\,\wt(Q_0)\le 256\cdot107$, so
$\wt(Q_0)=240$ and $\W_{Q_0}(0)=2B\ne0$. Since $e\in\Rad(Q_0)$, $Q_0(a\oplus e)=Q_0(a)\oplus\varepsilon_0$,
and $\varepsilon_0=1$ would make $Q_0$ balanced. So $Q_0(1,a'')=Q_0(0,a'')$.

\emph{(b) $m_1=0$.} Let $S^*=S(\varepsilon,\cdot)$, so $|\mathrm{supp}\,S^*|=m_1$. On $H_0$,
$\W_f(0,c,a'')=2(\W_{S(0,\cdot)}\pm\W_{S(1,\cdot)})(a'')\equiv 8Q_0(c,a'')\pmod{16}$, and by (a) the
right side does not depend on $c$. Adding and subtracting gives $4\mid\W_{S^*}$. Next,
\eqref{eq:star} gives $4\delta+64m_2\ge 800$, so $m_2\ge 11$ and $m_1=\delta-2m_2\le2$. Then
$|\W_{S^*}|\le m_1\le2$ and $4\mid\W_{S^*}$ force $\W_{S^*}=0$, and Parseval gives $m_1=0$.
Now $m_2=11$ would need $\delta=24=22+m_1$, so $m_2=12$ and $\delta=24$.

\emph{(c) Contradiction.} With $\delta=24$, \eqref{eq:delta} and hence \eqref{eq:L} are
equalities. So $s\in\{36,-28\}$ on all of $H_1$, i.e.\ $16+u\pm z\in\{18,-14\}$ for both signs.
Adding the two gives $u\in\{2,-14,-30\}$, so $u\equiv2\pmod4$. But $m_1=0$ means $E^*$ equals
$-2\sg(b)$ on $12$ points and $0$ elsewhere, so $\W_{E^*}$ is twice a sum of $12$ signs, and
$4\mid u$. \qed

\section{The case $n=8$}\label{sec:eight}

Here $B=8$, $t\in\{11,12\}$ and, by \eqref{eq:delta}, $\delta\le12$.

\emph{(a) $m_1\le2\Rightarrow m_1=0$.} By \eqref{eq:H0}, $48\,\wt(Q_0)\le 64\cdot47$, so
$\wt(Q_0)=56$. The argument of Section~\ref{sec:ten}(a),(b) then gives $4\mid\W_{S^*}$, and
hence $m_1=0$ whenever $m_1\le2$.

\emph{(b) The five local patterns.} On $H_1$, $s\equiv4\pmod{16}$ and $|s|\le20$, so
$s\in\{-12,4,20\}$, i.e.\ $u\pm z\in\{-14,-6,2\}$. Since $|u|\le\sum|E^*|=\delta\le12$,
\[
  (u,|z|)\in\{(-10,4),\,(-6,0),\,(-6,8),\,(-2,4),\,(2,0)\}.
\]
In particular $u\equiv2\pmod 4$ everywhere. If $m_1=0$ then
$\W_{E^*}\equiv 2m_2\pmod4$, so $m_2$ is odd.

\emph{(c) A second averaged inequality.} Each of the five pairs satisfies
$3z^2+5u^2+44u\le108$, with equality exactly at $(-10,4),(-6,8),(2,0)$. Summing with the
identities of Step~9 gives
\begin{equation}\label{eq:farkas}
  12t+5m_1+18m_2\le 216 .
\end{equation}

\emph{(d) The parameters are forced.} If $t=12$, then $\delta\le8$ and \eqref{eq:star} gives
$m_2\ge4$, so $(m_1,m_2)=(0,4)$, contradicting (b). So $t=11$. Then \eqref{eq:star} reads
$4m_1+40m_2\ge144$ and \eqref{eq:farkas} reads $5m_1+18m_2\le84$. This gives $m_2\le4$. Also
$m_2\le2$ would need $m_1\ge16>\delta$, and $m_2=4$ gives $m_1\le2$, hence $m_1=0$ and $m_2$
even, contradicting (b). So $m_2=3$, $m_1\ge6$, and with $\delta\le12$ we get
$(t,m_1,m_2)=(11,6,3)$. Both \eqref{eq:star} and \eqref{eq:farkas} are then equalities.

\emph{(e) Two values of $u$.} Equality in \eqref{eq:star} means $z^2=(2-u)^2$ at every $a''$,
and equality in \eqref{eq:farkas} means every pair lies in $\{(-10,4),(-6,8),(2,0)\}$.
Together, $u\in\{-6,2\}$ everywhere.

\emph{(f) A Boolean function with three large coefficients.} Since
$\W_{E^*}+2\tau=\tau(u+2)\in\{\pm4\}$, there is $g:\F^6\to\F$ with
$4\sg(g)=\W_{E^*}+2\tau$. By inversion,
\[
  \W_g=16E^*+4\,\sg\circ b ,
\]
so $|\W_g|=28$ at the $m_2=3$ points where $D^*=-\sg\circ b$, and $|\W_g|\le12$ elsewhere.

\emph{(g) Contradiction.} Let $x_1,x_2,x_3$ be those three points and $x_4=x_1\oplus x_2\oplus x_3$,
which is a fourth point, so $|\W_g(x_4)|\le12$. They form a coset $x_1+V$ of a
$2$-dimensional space $V$. For every character $\chi$ of $V$,
\[
  \Bigl|\sum_{y\in V}\chi(y)\W_g(x_1\oplus y)\Bigr|=|V|\,\Bigl|\sum_{x\in c+V^\perp}\sg(g(x))\sg(x_1\cdot x)\Bigr|\le 4\cdot16=64 ,
\]
where $c$ is the vector that defines $\chi$. Choosing $\chi$ so that the terms at $x_2,x_3$ have the
same sign as the term at $x_1$ gives a sum of absolute value at least $3\cdot28-12=72>64$. \qed

\section{The Lean formalization}\label{sec:lean}

The release contains two files for Lean~4.15.0 \cite{Lean4}. Neither imports anything, not even
Mathlib.
\begin{itemize}[leftmargin=1.5em]
\item \texttt{EvenBalanced8Bound.lean} (4{,}589 lines) proves Theorem~\ref{thm:main} for all
$m\ge4$ and Corollary~\ref{cor}.
\item \texttt{EvenBalancedBound.lean} (3{,}413 lines) is an earlier, smaller file for $m\ge5$.
Its content is contained in the first file. It is included because it checks faster and does
not use the $n=8$ argument.
\end{itemize}

\noindent\textbf{Definitions.} The statements use only elementary definitions, which a reader
can check against the usual ones:
\begin{verbatim}
abbrev Vec (n : Nat) := Fin n → Bool
def isBalanced (n : Nat) (F : Vec n → Bool) : Prop :=
  2 * hammingWeight F = (2 : Int) ^ n
def hasNonlinearityAtLeast (n : Nat) (F : Vec n → Bool) (d : Int) : Prop :=
  forall (a : Vec n) (c : Bool), d <= hammingDist F (affine a c)
\end{verbatim}
Here \texttt{affine a c x} is $(a\cdot x)\oplus c$, and \texttt{hammingWeight} and
\texttt{hammingDist} count over all of $\F^n$.

\noindent\textbf{Main statement.}
\begin{verbatim}
theorem even_balanced_nl_le_from8 (m : Nat) (hm : 4 <= m)
    (F : Vec (2 * m) → Bool) (hbal : isBalanced (2 * m) F) :
    ¬ hasNonlinearityAtLeast (2 * m) F
        ((2 : Int) ^ (2 * m - 1) - (2 : Int) ^ (m - 1) - 3)
\end{verbatim}
That is, no balanced $F$ has nonlinearity $\ge 2^{2m-1}-2^{m-1}-3$, which is
Theorem~\ref{thm:main}. Specializations include \texttt{balanced\_8\_nl\_le\_116},
\texttt{balanced\_10\_nl\_le\_492}, and S-box and permutation versions. The first file also
re-proves the known divisibility theorem of Sarkar and Maitra \cite{SM00} for resilient
functions. We do not claim that result as new.

\noindent\textbf{Availability.} The paper, the Lean files and the checking scripts are archived at
\url{https://doi.org/10.5281/zenodo.24546078}.

\noindent\textbf{Checking.} The release also contains a short file \texttt{ChallengeDefs.lean}
with the definitions above and a file \texttt{Challenge.lean} with the eight main statements, so a
reader only has to check these two short files. \texttt{verify.sh} checks the SHA-256 sums of all
Lean files and confirms that no \texttt{sorry}, \texttt{axiom} declaration or \texttt{native\_decide}
occurs in the code; the two proof files import nothing. It compiles everything, checks that
\texttt{Solution.lean} proves exactly the statements of \texttt{Challenge.lean} over identical
definitions, and prints the axioms of the main
theorems; the only ones are \texttt{propext}, \texttt{Classical.choice} and
\texttt{Quot.sound}. On an Apple-silicon laptop the two files compile in about $35$\,s
and $80$\,s, using at most about $1.2$ and $1.8$\,GB of memory. The proof follows
Sections~\ref{sec:main}--\ref{sec:eight} step by step. Its main lemma is the private theorem
\texttt{even\_balanced\_exclusion\_aux}.

\section{Remarks and questions}

\begin{enumerate}[leftmargin=1.5em]
\item For $n\ge12$ the final inequality of Section~\ref{sec:main} has slack. For instance,
$B=32$ gives $59\le t\le 58$. This suggests that $|\W_f|\le 2B+8$ might also be excluded for large
$m$, which would give $\nl\le 2^{2m-1}-2^{m-1}-6$. We have not proved this.
\item Is $\nl_{\mathrm{bal}}(12)=2010$, as Dobbertin's conjecture predicts? This would follow by
excluding $2012$, i.e.\ excluding $|\W_f|\le 72$ for balanced $f$ on $12$ variables.
\end{enumerate}

\section*{Acknowledgements}
The mathematical argument was developed by the author and the AI system Gemini Neon (Google)
working together interactively: the author contributed ideas, chose the directions to pursue and
took part in the reasoning at each step. Gemini Neon wrote the Lean code. Claude (Anthropic)
assisted with drafting this manuscript and preparing the release. Every proof was checked by the
Lean kernel, and the author takes full responsibility for the content of this paper.

\begin{thebibliography}{9}
\bibitem{Carlet21} C.~Carlet, \emph{Boolean Functions for Cryptography and Coding Theory},
Cambridge University Press, 2021.
\bibitem{Dob94} H.~Dobbertin, Construction of bent functions and balanced Boolean functions
with high nonlinearity, in: \emph{Fast Software Encryption (FSE 1994)}, LNCS 1008, Springer, 1995,
pp.~61--74.
\bibitem{LXW26} Y.~Lou, Y.~Xu and Q.~Wang, The maximum nonlinearity of 8-variable balanced
Boolean functions is 116, arXiv:2610.17897, 2026.
\bibitem{Lean4} L.~de~Moura and S.~Ullrich, The Lean~4 theorem prover and programming language,
in: \emph{CADE-28}, LNCS 12699, Springer, 2021, pp.~625--635.
\bibitem{SM00} P.~Sarkar and S.~Maitra, Nonlinearity bounds and constructions of resilient
Boolean functions, in: \emph{CRYPTO 2000}, LNCS 1880, Springer, 2000, pp.~515--532.
\bibitem{SZZ93} J.~Seberry, X.-M.~Zhang and Y.~Zheng, Nonlinearly balanced Boolean functions and
their propagation characteristics, in: \emph{CRYPTO '93}, LNCS 773, Springer, 1994, pp.~49--60.
\bibitem{XW26} Y.~Xu and Q.~Wang, The maximum possible nonlinearity for a balanced Boolean
function in 10 variables is 492, IACR Cryptology ePrint Archive, Report 2026/2797, 2026.
\bibitem{YHZZW26} Y.~Yuan, Y.~Hu, Y.~Zhang, L.~Zhang and W.~Wu, Theoretical open problems in
symmetric cryptography: verifiable LLM-guided analysis, IACR Cryptology ePrint Archive,
Report 2026/1687, 2026.
\end{thebibliography}

\end{document}
